% =============================================================================
% CTSpinoPelvic1K — Visual Abstract / dataset-construction flow.
%
% Compile standalone (cropped figure):   pdflatex visual_abstract.tex
% (this is kept in sync with fig:visual_abstract in full-paper-template.tex).
%
% The "book-page" stacks are PLACEHOLDERS. To drop in a real 2-D CT-with-mask
% overlay, redefine the matching slot before \begin{document}, e.g.
%     \renewcommand{\rSpine}{\includegraphics[width=1.9cm,height=2.4cm]{renders/spine.png}}
% =============================================================================
\documentclass[border=12pt]{standalone}
\usepackage{graphicx}
\usepackage{tikz}
\usetikzlibrary{arrows.meta, calc, positioning}

\definecolor{cspine}{RGB}{120,42,0}
\definecolor{cpelv}{RGB}{80,22,100}
\definecolor{cfuse}{RGB}{25,100,180}
\definecolor{ctrain}{RGB}{18,78,40}
\definecolor{cpseudo}{RGB}{175,95,0}
\definecolor{crib}{RGB}{150,30,30}
\definecolor{cgray}{RGB}{90,95,105}

\tikzset{
  pback/.style={draw=#1!55!black, fill=#1!18, rounded corners=2pt,
                minimum width=2.0cm, minimum height=2.5cm, line width=0.4pt},
  pfront/.style={draw=#1!70!black, fill=#1!7, rounded corners=2pt,
                 minimum width=2.0cm, minimum height=2.5cm, line width=0.8pt,
                 inner sep=1.4pt},
  flow/.style={-{Stealth[length=2.6mm]}, line width=1.0pt, draw=cgray},
  hdr/.style={font=\sffamily\bfseries\small, text=#1!55!black, align=center},
  lbl/.style={font=\sffamily\small, align=center, text=#1!50!black},
  cnt/.style={font=\sffamily\bfseries\normalsize, text=black},
  procbox/.style={draw=#1!70!black, fill=#1!9, rounded corners=5pt,
                  line width=1.0pt, align=center, font=\sffamily, inner sep=8pt},
}

% Three offset "pages"; #5 = page content (a placeholder now / \includegraphics later).
\newcommand{\pagestack}[5]{%
  \node[pback=#4] at ($(#2,#3)+(0.24,0.24)$) {};
  \node[pback=#4] at ($(#2,#3)+(0.12,0.12)$) {};
  \node[pfront=#4] (#1) at (#2,#3) {#5};
}
% Placeholder shown until a real overlay is supplied.
\newcommand{\vph}[1]{\parbox[c][2.2cm][c]{1.75cm}{\centering
  \scriptsize\itshape\color{cgray}#1\\[3pt]\tiny[2-D render]}}
% ── RENDER SLOTS — redefine to insert a real CT-with-mask overlay ──
\newcommand{\rSpine}     {\vph{spine GT}}
\newcommand{\rPelvis}    {\vph{pelvic GT}}
\newcommand{\rFused}     {\vph{spine+pelvis}}
\newcommand{\rSpineOnly} {\vph{spine}}
\newcommand{\rPelvicOnly}{\vph{pelvis}}
\newcommand{\rDense}     {\vph{spine+pelvis}}
\newcommand{\rRibs}      {\vph{+ ribs}}

\begin{document}
\begin{tikzpicture}[x=1cm, y=1cm]
\useasboundingbox (-3.9,-5.7) rectangle (20.1,6.2);

% ── phase headers (kept well above the stacks) ──
\node[hdr=cgray]   at (0,    5.5) {\textsc{Source masks}\\[-1pt]\footnotesize two TCIA datasets};
\node[hdr=cgray]   at (5,    5.5) {\textsc{Patient-anchored}\\[-1pt]\footnotesize match (bone-HU)};
\node[hdr=ctrain]  at (9.6,  5.5) {\textsc{Partial-annotation}\\[-1pt]\footnotesize training};
\node[hdr=cpseudo] at (13.8, 5.5) {\textsc{Pseudolabel}\\[-1pt]\footnotesize complete the pelvis};
\node[hdr=crib]    at (17.8, 5.5) {\textsc{+ Ribs}\\[-1pt]\footnotesize TotalSegmentator, v3};

% ── (1) source mask stacks (labels to the LEFT, clear of the merge) ──
\pagestack{spine}{0}{2.5}{cspine}{\rSpine}
\node[lbl=cspine, anchor=east, align=right] at (-1.2,2.5)
  {CTSpine1K\\spine GT\\\textbf{784 masks}};
\pagestack{pelv}{0}{-2.5}{cpelv}{\rPelvis}
\node[lbl=cpelv, anchor=east, align=right] at (-1.2,-2.5)
  {CTPelvic1K\\pelvic GT\\\textbf{714 masks}};

% ── (2) match: } brace -> fused; -> arrows -> separate ──
\pagestack{sonly}{5}{3.4}{cspine}{\rSpineOnly}
\pagestack{fused}{5}{0}{cfuse}{\rFused}
\pagestack{ponly}{5}{-3.4}{cpelv}{\rPelvicOnly}
\node[cnt, anchor=west, fill=white, inner sep=1.5pt] at (6.15,3.4) {Spine-only \textbf{440}};
\node[cnt, anchor=west, fill=white, inner sep=1.5pt] at (6.15,0.85) {Fused \textbf{342}};
\node[cnt, anchor=west, fill=white, inner sep=1.5pt] at (6.15,-3.4) {Pelvic-only \textbf{371}};

% the } merge, drawn as an explicit path (spike points at Fused)
\draw[line width=0.8pt, draw=cfuse!55!black] (spine.east) -- (2.4,2.5);
\draw[line width=0.8pt, draw=cfuse!55!black] (pelv.east)  -- (2.4,-2.5);
\draw[line width=1.3pt, draw=cfuse!75!black, rounded corners=2pt]
  (2.4,2.5) -- (2.7,2.5) -- (2.7,0.35) -- (3.05,0) -- (2.7,-0.35) -- (2.7,-2.5) -- (2.4,-2.5);
\draw[flow, draw=cfuse!75!black] (3.05,0) -- (fused.west);
\node[lbl=cfuse, font=\sffamily\scriptsize, fill=white, inner sep=1pt] at (3.55,1.25) {same scan};
% separate cases: straight arrows
\draw[flow] (spine.east) to[out=0,in=180] (sonly.west);
\draw[flow] (pelv.east)  to[out=0,in=180] (ponly.west);
\node[lbl=cgray, font=\sffamily\scriptsize, fill=white, inner sep=1pt] at (2.55,3.75) {different scans};

% ── (3) training box ──
\node[procbox=ctrain, minimum width=2.4cm, minimum height=5.0cm] (train) at (9.6,0)
  {\bfseries nnU-Net v2\\[2pt]ignore-label\\training\\[7pt]
   \footnotesize\textbf{all 1{,}153}\\\footnotesize cases};
\draw[flow] (fused.east) -- (train.west);
\draw[flow] (sonly.east) to[out=0,in=155] ($(train.west)+(0,1.8)$);
\draw[flow] (ponly.east) to[out=0,in=205] ($(train.west)+(0,-1.8)$);

% ── (4) pseudolabel completion ──
\pagestack{dense}{13.8}{0}{cpseudo}{\rDense}
\node[lbl=cpseudo, align=center] at (13.8,-1.95) {dense spine\,+\,pelvis\\\textbf{782}};
\draw[flow, draw=cpseudo!75!black] (train.east) -- (dense.west)
  node[midway, above=2pt, font=\sffamily\scriptsize, text=cpseudo!60!black, align=center]
  {fused 342\\+ spine-only 440\\\textbf{= 782}};
\node[lbl=cgray, font=\sffamily\scriptsize, align=center] at (11.7,-3.2)
  {pelvic-only 371\\held out};

% ── (5) ribs (TotalSegmentator) ──
\pagestack{ribs}{17.8}{0}{crib}{\rRibs}
\node[lbl=crib, align=center] at (17.8,-1.95) {spine\,+\,pelvis\,+\,ribs\\\textbf{782}};
\draw[flow, draw=crib!75!black] (dense.east) -- (ribs.west)
  node[midway, above=2pt, font=\sffamily\scriptsize, text=crib!60!black, align=center]
  {TotalSeg.\\ribs};

\end{tikzpicture}
\end{document}
